It is based on the notion of rigged conf\/igurations~\mbox{\cite{KKR,KR}}, and is described in the following way.
Let $m_j$ be the number of solitons with amplitude~$j$.
In what follows $m_j$ are supposed to be positive for $1 \leq j \leq N$ and to be zero for $j > N$ for simplicity.
One can think of $\m=(m_j)_{1 \leq j \leq N}$ as a collection of f\/ixed values of the f\/irst integrals.
We denote by $\J (\m)$ the level set specif\/ied by $\m$
and by $\Omega(\m)$
its volume.
Then we have \cite{KTT}
\begin{equation}
\Omega(\m) =
(\det F)\prod_{1 \leq j \leq N} \frac{1}{m_j}
\binom{p_j + m_j - 1}{m_j - 1}.\label{eq:oct10_1}
\end{equation}
Here $p_j$ are positive integers and $F$ is an $N \times N$ matrix with integer entries,
which are explicitly expressed in terms of $\m$ and the system size $L$.

{It can be shown that}
if there is no solitons of common amplitudes then
the level set $\J (\m)$ becomes a torus $F \Z^N \backslash \Z^N $.\footnote{More precisely, it is not a torus but a set of all integer points on the torus $F \Z^N \backslash \R^N $.
For the sake of simplicity we call such a set a torus.
Besides, in this paper we write all the quotient sets as left quotient ones.}
This fact is suggested by
the volume formula (\ref{eq:oct10_1}) that reduces to $\Omega(\m) =\det F$ when $m_j=1$ for every $j$.
However this simple picture fails when there are multiple solitons of common amplitudes.
{If so,} the right hand side of (\ref{eq:oct10_1}) can not be regarded as the volume of a torus any more.
{An expression for} the level set in this generalized case is given by
\begin{equation}\label{eq:jv}
\J (\m) =  A \Z^{m_1+\cdots + m_{N}} \backslash ({\mathcal I}_{m_1}\times
\cdots \times {\mathcal I}_{m_N}),
\end{equation}
where ${\mathcal I}_n = {\mathfrak S}_n \backslash (\Z^n-\Delta_n)$ is
the $n$ dimensional lattice without the diagonal
